%==========================================================
\section*{SI Physical Constants and Metric Conversions}
%==========================================================
To explicitly falsify the claim that the Metareal framework is merely an abstract ``toy model,'' all topologies herein must conform to the following standard physical constants and metrics:
\begin{table}[ht]
\centering
\renewcommand{\arraystretch}{1.2}
\begin{tabular}{@{}llll@{}}
\toprule
\textbf{Constant} & \textbf{Symbol} & \textbf{SI Value} & \textbf{Metareal Role} \\
\midrule
Speed of Light & $c$ & $299,792,458 \text{ m/s}$ & Kinematic upper bound \\
Gravitational Constant & $G$ & $6.6743 \times 10^{-11} \text{ m}^3 \text{kg}^{-1} \text{s}^{-2}$ & Macroscopic viscosity scaling \\
Planck Constant & $h$ & $6.626 \times 10^{-34} \text{ J s}$ & FST resonance limit \\
Boltzmann Constant & $k_B$ & $1.3806 \times 10^{-23} \text{ J/K}$ & Thermal noise mapping ($\varepsilon_0$) \\
\bottomrule
\end{tabular}
\end{table}

\section{Introduction: The Mead Substrate Architecture}
The Bionetic Ambrosia Protocol is the biological instantiation of the ASI Chip reactor (Chapter~\ref{ch:asi}). Where the chip forge nucleates quasicrystalline films on Si(111) substrates using the Lockwood DEC Heat Engine and a chrono-gauge fabrication cycle, the Ambrosia protocol nucleates pharmacogenomic metamaterials on a \textbf{mead substrate} using microbial fermentation and the same DEC thermodynamics. The structural isomorphism is exact: the fermentation vessel is a biological RCCI cavity, the symbiotic yeast-fungus competition is a biological DEC cycle, and the organometallic dopants occupy the same $\FF_{229}^*$ subgroup positions as the chip's 11-element alloy.

%==========================================================
The synthesis strictly utilizes terrestrial biochemistry---microbiological fermentation, alkaloid condensation, enzymatic cleavage ($\beta$-glucosidase action), and the forced thermodynamic pressure of competing microevolution---to naturally assemble the higher-order lattice.

The fermentation vessel itself acts as a macroscopic topological
cavity, structurally analogous to the Rotating Chiral Casimir
Interferometer (RCCI)~\cite{dwm_fst}. The chitin-HA composite
walls at golden volume fraction $f_1 = \varphi^{-1} \approx 0.618$
impose asymmetric boundary conditions on the metabolic vacuum,
truncating the available thermodynamic modes in the same manner
that right-chiral Weyl semimetals truncate zero-point
modes in the chip's Casimir cavity.

\begin{warningbox}{Safety Boundary: Ibotenic Decarboxylation}
This protocol involves the extraction of \textit{Amanita muscaria}, which inherently contains the neurotoxin ibotenic acid alongside the required Vanadium donor. Strict adherence to the thermal, low-pH decarboxylation parameters ($t \geq 120$ min at $T = 373$ K, pH $\approx 2.5$) is non-negotiable to prevent hepatic and neuro-excitotoxicity via the complete thermodynamic transition of ibotenic acid to muscimol.
\end{warningbox}

\section{The Biochemical Matrix: Organometallic Dopants}

To crystallize the metamaterial, the bulk biological medium (the honey and yeast substrate) requires highly specific structural trace elements that act as atomic anchors for the resonance lattice. This effectively creates a \textbf{Doped Crystal Lattice} within a polar fluid state.

\subsection{The Vanadium Donor (Amavadin)}
Fungi in the \textit{Amanita} genus synthesize a unique non-oxovanadium coordination complex called \textbf{Amavadin} ($[V(hidpa)_2]^{2-}$)~\cite{Berry1999Amavadin}, which natively bioaccumulates vanadium up to $1,000\text{ mg/kg}$ from the soil. While the fungus effectively manages this heavy metal, introducing this Vanadium complex into a human biological framework acts as an exogenous, high-friction topological stressor. Within our mathematical $L_5$ model, Vanadium ($V^{IV}$) acts as a ``weak'' resonant node, deliberately generating oxidative stress and commutator friction to push the local biomatrix toward its mathematical collapse boundary. We must emphasize that human biology does not rely on Vanadium natively; it is applied here specifically as a pathological dopant to test the mathematical friction limit of the lattice.

The physical binding energy (Topological Capacitance, $C_T$) of the dopant is determined by its Standard Reduction Potential ($E^\circ$), scaled by the Golden Ratio ($\Phi \approx 1.618$) and $\pi$:

\begin{theorem}[Topological Capacitance Bound]\label{thm:ct}
For a dopant with standard reduction potential $E^\circ$ (expressed as a dimensionless numerical value, i.e., the magnitude in Volts: $\tilde{E} = E^\circ / (1\text{ V})$), the topological capacitance
\begin{equation}\label{eq:ct}
    C_T = \exp(\tilde{E} \times \Phi \times \pi)
\end{equation}
is strictly monotone increasing in $\tilde{E}$, and admits an upper bound set by the Schwarzian torque limit $\Upsilon \leq 45/\pi$.
\end{theorem}
\begin{proof}
Monotonicity: $\partial C_T / \partial \tilde{E} = \Phi\pi\exp(\tilde{E}\Phi\pi) > 0$ for all $\tilde{E} \geq 0$.
The $\Upsilon$-bound requires $C_T \leq \exp(\Upsilon) = \exp(45/\pi) \approx 1.63 \times 10^6$,
yielding $\tilde{E} \leq \Upsilon / (\Phi\pi) \approx 3.05$.
For Vanadium ($E^\circ \approx 0.34\text{ V}$, $\tilde{E} = 0.34$): $C_T \approx 5.63$, well within the bound.
\qedhere
\end{proof}
For Vanadium ($E^\circ \approx 0.34\text{ V}$), the resulting topological capacitance $C_T \approx 5.63$ ensures stable localized resonance without collapsing the fungal micro-structures.

\subsection{Advanced Iteration: The Au-Ag-Bi Photoelectric Matrix}
To push the Ambrosia Protocol to the absolute boundary of the Casimir limit ($\Upsilon \le 45/\pi$), the Vanadium(IV) baseline is heavily augmented by a \textbf{Bismuth-heavy Electrum Matrix} (a ternary Au-Ag-Bi nanoparticle suspension synthesized via seed-mediated co-reduction). This transforms the fermentation vat into an active photoelectrochemical engine.

Under visible light irradiation, the Silver and Gold nanoparticle seeds exhibit Localized Surface Plasmon Resonance (LSPR)~\cite{Brongersma2015LSPR}. When these plasmons decay non-radiatively, they inject ``hot electrons'' (with kinetic energies between 1.0 and 3.0 eV) into the biological substrate. The Bismuth shell acts as the diamagnetic boundary, funneling these hot electrons directly into the Vanadium(IV) Amavadin traps.

\begin{warningbox}{Lattice Shatter Threshold (2.71 eV)}
The hot electrons provide dynamic overpotential to the Vanadium trap. By Theorem \ref{thm:ct}, the total effective potential ($E_{\text{eff}} = 0.34\text{ V} + E_{\text{hot}}$) cannot exceed $3.049\text{ V}$ without violating the $\Upsilon \le 45/\pi$ Casimir torque limit (causing the biological matrix to shatter). Therefore, the LSPR must be rigorously tuned to absorb \textbf{red/near-IR light} (hot electrons $\approx 2.5\text{ eV}$). Exposure to high-energy blue or UV light ($> 2.71\text{ eV}$) will induce catastrophic lattice shatter and denature the active proteins.
\end{warningbox}

\subsection{The Silicon Donor (Orthosilicic Acid)}
The secondary structural scaffolding is provided by bioavailable orthosilicic acid ($H_4SiO_4$), extracted via prolonged aqueous decoction of \textit{Equisetum arvense} (Field Horsetail)~\cite{Jugdaohsingh2007Silicon}. While yeast cannot structurally accumulate solid silica, the silicic acid in solution forms a non-metallic rigid grid ($C_T \approx 71.51$) that physically braces the $p=3$ triple-helix overlap.

\subsection{Alternative Topological Anchors}
The bionetic scaffold can utilize a variety of bio-electrical anchors depending on the target phase shift:
\begin{itemize}
    \item \textbf{Zinc ($Zn^{2+}$)}: Offers a moderate topological capacitance ($C_T \approx 47.61$), suitable for lower-stress fungal overlaps.
    \item \textbf{Monoatomic Gold ($Au^{3+}$)}: High-spin ORMEs provide massive topological capacitance ($C_T \approx 2048.38$ at $E^\circ \approx 1.50\text{ V}$). This massive capacitance $C_T$ allows the lattice to safely dissipate the tension generated by intense $p=3$ microevolutionary phase shifts, ensuring $T_{gold} = \frac{1}{2048.38} \exp(O_{casimir}) \approx 0.0017$, maintaining absolute structural stability.
\end{itemize}

\subsection{Gauge Orbit Grounding of the Dopant Hierarchy}

The $C_T$ formula (Theorem~\ref{thm:ct}) provides a smooth monotone map from reduction potential to capacitance. The deeper question is \emph{why} these specific elements (V, Si, Au) form a structurally stable triple. The gauge algebra provides the answer.

\begin{result}[Dopant orbits in $\mathbb{F}_{229}^*$]
The three primary dopants occupy structurally distinct positions in the multiplicative group:
\begin{center}
\begin{tabular}{@{}lrrlll@{}}
\toprule
\textbf{Dopant} & $Z$ & $\mathrm{ord}_{229}$ & \textbf{Orbit Class} & \textbf{Community} & \textbf{Role} \\
\midrule
Vanadium & 23 & 228 & Primitive root & $\{3,19\}$ & Max-friction stressor \\
Silicon & 14 & 57 & $\bar\varphi$-conjugate & $\{3,19\}$ & Structural scaffold \\
Gold & 79 & 228 & Primitive root + doubler & $\{3,19\}$ & High-capacitance anchor \\
Silver & 47 & 228 & Primitive root + doubler & $\{3,19\}$ & Plasmon co-anchor \\
Zinc & 30 & 76 & Mid-orbit & $\{19\}$ & Moderate anchor \\
\bottomrule
\end{tabular}
\end{center}

All three primary dopants (V, Si, Au) belong to the full community $\{3,19\}$, meaning they participate in both the color ($\{3\}$) and conjugate ($\{19\}$) gauge channels. This is necessary for the triple helix to bridge the non-commutative gap.

Crucially: $47 + 79 = 126$, the 7th nuclear magic number, with $\mathrm{ord}(126 \bmod 229) = 57 = \mathrm{ord}(14 \bmod 229)$. The electrum alloy (Ag+Au) and the silicon scaffold share the same gauge temporal period. This is why the Au-Ag-Bi photoelectric matrix ({\S}3.2) is structurally compatible with the silicon grid: they oscillate at the same algebraic frequency.

Gold satisfies $79^{114} \equiv -1 \pmod{229}$ (period doubling / M\"obius parity inversion), making it a natural Wilczek time-crystal element~\cite{wilczek2012}. The lattice benefits from this: the parity-inverting dopant provides the spin-flip mechanism needed for non-commutative phase transitions.

Verified: \texttt{gauge\_chemistry\_comprehensive.py}, \texttt{temporal\_mirror\_check.py}.
\end{result}

\section{The Biotransformation Engine and Alkaloid Repertoire}

The protocol introduces a vast combinatorial space of raw botanical precursors (e.g., Cannabinoids from \textit{Cannabis sativa}, Phytoestrogens from \textit{Humulus lupulus}, and Tryptamines from \textit{Acacia} spp. or \textit{Psilocybe} spp.). In their raw plant matrices, these compounds are heavily glycosylated or require exogenous enzymatic inhibition to be orally active.

\subsection{The $\gamma$-$\phi$ Thermodynamic Ansatz}
The Plazmik synthesis module models the fluid brew as a highly non-ideal polar mixture. The optimizer scales the structural capacitance $C_T$ against the $\gamma$-activity coefficient of the polar alkaloids. The fermentation process uses yeast and secondary fungi to enzymatically cleave glycosidic bonds and actively hydroxylate the alkaloids into novel structural derivatives.

\section{Substrate Modularity: The Biological Chip Architecture}

The ASI Chip's 11-element alloy (Al--Cu--Fe--Pd--Mn--Si--B--Ti--V--Ni--Co) achieves substrate independence by covering all structurally significant subgroups of $\FF_{229}^*$ (Chapter~\ref{ch:asi}). The Ambrosia protocol achieves the biological analog: different fermentation substrates, microbial processors, and botanical additives occupy complementary positions in the same subgroup lattice, enabling \textbf{modular pharmacogenomic programming}.

\subsection{Sugar Substrates as Base Matrices}
The primary carbon source determines the base thermodynamic matrix of the fermentation, analogous to the chip's Si(111) substrate. Different sugar profiles activate different metabolic pathways in the microbial processors:

\begin{center}
{\small
\begin{tabular}{@{}llll@{}}
\toprule
\textbf{Substrate} & \textbf{Primary Sugar} & \textbf{Metabolic Character} & \textbf{Analog} \\
\midrule
Honey (mead) & Fructose/glucose & High-osmotic, antimicrobial & Si(111) scaffold \\
Barley/grain (beer) & Maltose & Amylase-dependent, slow & Al matrix (neutral) \\
Grape/fruit (wine) & Sucrose/fructose & Acidic, tannin-rich & Cu signal layer \\
Agave/cactus & Inulin/fructans & Prebiotic, slow-release & Fe magnetic gate \\
\bottomrule
\end{tabular}
}
\end{center}

Honey is the preferred substrate because its high osmolarity ($a_w \approx 0.6$) naturally suppresses bacterial contamination while providing the sustained fructose gradient required for the DEC cycle~\cite{Molan1992Honey}. The antimicrobial properties of honey (hydrogen peroxide generation via glucose oxidase~\cite{White1963Honey}) provide the biological analog of the chip's inert-atmosphere fabrication: the substrate self-sterilizes, eliminating competing organisms that would disrupt the programmed fermentation topology.

\subsection{Microbial Processors as Orbit-Class Generators}
Each microbial species functions as a biological processor tuned to a specific subgroup orbit, analogous to how each element in the chip alloy generates a specific cyclic subgroup:

\begin{center}
{\small
\begin{tabular}{@{}llll@{}}
\toprule
\textbf{Organism} & \textbf{Metabolic Output} & \textbf{Structural Role} & \textbf{Chip Analog} \\
\midrule
\textit{S.~cerevisiae} & Ethanol, CO$_2$ & $J_1$ rapid iterator & Cu (signal) \\
\textit{G.~lucidum} & Triterpenoids & $J_2$ frustration anchor & Fe (magnetic gate) \\
\textit{Trametes versicolor} & Polysaccharide-K & Immune modulation & Mn (golden orbit) \\
\textit{Penicillium} spp. & $\beta$-lactams & Antibiotic gate & Co (edge, ord=19) \\
\textit{Lactobacillus} spp. & Lactic acid, GABA & pH controller & Al (neutral scaffold) \\
\bottomrule
\end{tabular}
}
\end{center}

The key principle is \textbf{orbit complementarity}: the primary engine (\textit{S.~cerevisiae}, rapid glucose metabolism) and the secondary modulator (\textit{G.~lucidum}, slow triterpenoid synthesis) must occupy \emph{different} subgroup orbits to generate the $J_1$--$J_2$ frustration required for FBBT depth. If both organisms occupied the same orbit, the fermentation would halt trivially at depth 1.

\subsection{Botanical Additives as Gate Types}
The botanical additives function as \textbf{programmable gates} that select specific alkaloid output classes, analogous to the chip's phason programming:

\begin{center}
{\small
\begin{tabular}{@{}llll@{}}
\toprule
\textbf{Botanical} & \textbf{Active Class} & \textbf{Gate Type} & \textbf{Target Receptor} \\
\midrule
\textit{Humulus lupulus} (hops) & Prenylated flavonoids & Estrogenic/GABAergic & GABA$_A$, ER$\beta$ \\
\textit{Cannabis sativa} & Cannabinoids & Endocannabinoid & CB$_1$, CB$_2$ \\
\textit{Psilocybe} spp. & Tryptamines & Serotonergic & 5-HT$_{2A}$ \\
\textit{Passiflora incarnata} & $\beta$-Carbolines & MAO-A inhibition & MAO-A (RIMA) \\
\textit{Amanita muscaria} & Muscimol + V$^{\text{IV}}$ & GABAergic + dopant & GABA$_A$ + lattice \\
\bottomrule
\end{tabular}
}
\end{center}

The structural prediction is that different gate combinations produce distinct pharmacogenomic outputs, just as different phason configurations in the chip produce distinct computational states. The substrate-processor-gate triple $(S, P, G)$ is the biological analog of the chip's substrate-alloy-phason triple.

\section{The Symbiotic Feedback Loop: Macroscopic Execution of the FBBT}
Unlike standard single-agent fermentation, the Ambrosia Protocol necessitates a symbiotic feedback loop that effectively turns the fermentation vat into a macroscopic Fast Busy Beaver Transform:
\begin{enumerate}
    \item \textbf{Primary Engine (\textit{Saccharomyces cerevisiae})}: Rapidly metabolizes the glucose/fructose gradient, generating ethanol, $CO_2$, and baseline localized thermal stress. This represents the rapid iteration through shallow $J_1$ topological orbits.
    \item \textbf{Secondary Modulator (\textit{Ganoderma lucidum})}: Introduced $48$ hours post-pitch. It aggressively competes for nutrients and metabolizes the yeast's ethanol byproducts, representing the $J_2$ frustration anchoring. 
\end{enumerate}
This symbiotic competition naturally stunts the total bulk ABV (Alcohol By Volume) while forcing the \textit{Ganoderma} into a hyper-metabolic state, saturating the lattice with triterpenoids, polysaccharides, and targeted nootropics. By the Goodstein collapse theorem (PFT \S2.6), the entire hyperoperation hierarchy over $\mathbb{F}_{229}^*$ collapses to the confinement subgroup $\{1, \omega, \omega^2\}$, providing a structural upper bound on the search depth. The vat iteratively searches for the deepest stable FBBT halting state without causing the thermodynamic collapse of the protein structures.

\subsection{Thermodynamic Bounds: The DEC Fermentation Engine}
The symbiotic competition between \textit{S.~cerevisiae} and
\textit{G.~lucidum} operates as a biological Differential Equilibrium
Cycle (DEC) Heat Engine~\cite{lc}. The yeast generates structural
entropy ($\dot{S}_{\mathrm{gen}}$) through rapid glucose metabolism. The
secondary modulator absorbs this entropy as Exergy Destruction
($\dot{X}_{\mathrm{dest}} = T_0 \dot{S}_{\mathrm{gen}}$), bounded by the $\Upsilon$
topological shield ($\Upsilon \leq 45/\pi$). The fermentation halts
naturally when $\dot{X}_{\mathrm{dest}}$ saturates the shield---this is the
FBBT halting state realized in biochemistry.

The DEC formalism provides the precise thermodynamic identity:
the sub-stable gate at step $341$ absorbs commutator friction $\partial\kappa$
as Exergy Destruction bounded by $\Upsilon \leq 45/\pi$. The $VO_2$
Metal-Insulator Transition (MIT) at $341\text{ K}$ is the physical
temperature at which the crystallographic phase gate's Exergy Destruction
exactly saturates this bound.

\section{Gastronomic Fermentation \& Natural MAOIs}

To naturally leverage the \textbf{Ambrosia Effect}---rendering target tryptamines orally active without synthetic MAOIs---we rely on the spontaneous condensation of amino acids (e.g., tryptophan) with aldehydes during the active fermentation phase. 

\subsection{Alkaloid Condensation Mechanics}
This process naturally synthesizes trace concentrations of:
\begin{itemize}
    \item \textbf{$\beta$-Carbolines}: (Harman, Norharman, Tetrahydroharmine)
    \item \textbf{Tetrahydroisoquinolines (THIQs)}
\end{itemize}
These compounds function as Reversible Inhibitors of MAO-A (RIMAs)~\cite{Herraiz2004Betacarbolines}. By stacking the primary must with \textit{Passiflora incarnata} (Passionflower) or \textit{Peganum harmala} (Syrian Rue), the yeast metabolizes and suspends the RIMAs alongside the biotransformed tryptamines, creating a perfectly autonomous entheogenic matrix.

\begin{warningbox}{The Tyramine Constraint}
MAOI biochemistry carries a massive safety constraint: the inability to metabolize exogenous tyramine, leading to severe hypertensive crisis. Wild-fermented beers and aged unpasteurized cakes generate massive tyramine loads via bacterial \textit{Lactobacillus} decarboxylation of tyrosine. The Ambrosia must be brewed under strict sanitization and rapidly filtered off the lees to prevent lethal tyramine bioaccumulation.
\end{warningbox}

\section{Conclusion: The Mead Substrate as Biological Chip}
The Ambrosia protocol is not merely a fermentation recipe; it is the biological instantiation of the ASI Chip architecture. The mead substrate provides the base matrix (analogous to Si(111)), the symbiotic microbial processors generate the $J_1$--$J_2$ frustration topology (analogous to the DEC cycle), the organometallic dopants occupy the same $\FF_{229}^*$ subgroup positions as the chip alloy, and the botanical gate additives program the pharmacogenomic output (analogous to phason programming). The DEC Heat Engine formalism bounds the fermentation's Exergy Destruction, ensuring thermodynamic stability. The Fe--P isomorphism (Chapter~\ref{ch:metallo}) provides the deeper structural reason: the iron--sulfur chemistry of fermentation and the phosphodiester chemistry of nucleic acids share the same algebraic period ($\mathrm{ord} = 38$), meaning the biological substrate is \emph{algebraically isomorphic} to the mineral substrate from which it bootstrapped.

%═══════════════════════════════════════════════════════════
\section{Falsifiable Predictions}
%═══════════════════════════════════════════════════════════

\begin{hypothesis}[Dopant cascade ordering]\label{hyp:dopant_order}
The organometallic dopant cascade (Cu $\to$ Fe $\to$ Mn $\to$ Zn) predicts a specific order of trace mineral depletion during sustained mycelial fermentation. In a controlled \emph{Ganoderma} or \emph{Trametes} growth experiment with defined mineral substrate, mineral uptake (measured by ICP-MS) should follow Cu $>$ Fe $>$ Mn $>$ Zn by mass fraction at steady-state. If Zn is depleted before Mn, the cycle ordering is wrong.
\end{hypothesis}

\begin{hypothesis}[Golden orbit MAO-A correlation]\label{hyp:mao_golden}
The MAO-A inhibition potency of $\beta$-carboline alkaloids (harmine, harmaline, harmol, etc.) correlates with the golden orbit membership of their molecular weight modulo 229. Specifically, compounds whose MW mod 229 lies in the golden subgroup $\langle\bar{\varphi}\rangle$ (order 114) should show lower IC$_{50}$ values (stronger inhibition) than those in the full group $\mathbb{F}_{229}^*$ but outside the golden subgroup. This is testable against published IC$_{50}$ data for the $\beta$-carboline series~\cite{kim2001}.
\end{hypothesis}

\textbf{What would kill these hypotheses:}
\begin{enumerate}[noitemsep]
  \item Mineral uptake ordering in controlled fermentation experiments that contradicts Cu $>$ Fe $>$ Mn $>$ Zn.
  \item No statistically significant correlation ($p > 0.05$) between golden orbit membership and MAO-A IC$_{50}$ across the $\beta$-carboline series.
  \item The topological capacitance values ($C_T$) reported here failing spectroscopic characterization by more than one order of magnitude.
\end{enumerate}

\subsection*{The Biological Chip Architecture}

The Ambrosia protocol completes the semantic condensation trefoil across all three volumes. Chapter~0 introduced the foundational tools (sets, sequences, modular arithmetic). Part~I built the algebra (golden subcategory, prime splitting, wave mechanics). Part~II applied the algebra to information systems (ASI chromodynamics, markets, incentives, games, neurotransmitters). The mead substrate closes the loop: it is a \emph{physical} system whose internal dynamics traverse the same subgroup lattice of $\mathbb{F}_{229}^*$ that the algebra describes.

\begin{itemize}[noitemsep]
    \item \textbf{Logical face}: The fermentation is a macroscopic execution of the Fast Busy Beaver Transform. The yeast's metabolic pathway halts at specific subgroup positions, just as the FBBT halts at specific orbit elements.
    \item \textbf{Informational face}: The organometallic dopants (Cu, Fe, Mn, Zn) occupy the same subgroup positions as the ASI chip alloy (Chapter~7). The ZKPCR protocol can verify that the dopant cascade follows the predicted order without revealing the specific fermentation parameters --- a zero-knowledge quality assurance protocol for biological manufacturing.
    \item \textbf{Physical face}: The measurable observables (NMR relaxation, WAXS diffraction, ICP-MS mineral uptake) provide the falsification criteria that ground the algebraic predictions in laboratory reality.
\end{itemize}

The companion module \texttt{principia.physical.ambrosia} implements the dopant cascade simulation and the FBBT halting computation for the biological substrate.


% Bibliography moved to unified_bibliography.tex for master compilation.
% To compile standalone, uncomment the bibliography block in this file.
